% Research fragment. Requires the source-coordinate definitions in
% order_ten_completion.tex, or equivalent definitions from the parent article.
% No external computational package is required by the complete replay.
\section{Complete optimal reconstruction at smaller tail order twelve}
\label{t12:section}

Research question R1 of the source \emph{Quartic Spectral Reconstruction}
explicitly asks for complete smaller-order-ten and smaller-order-twelve
families. The preceding section and the theorem below resolve those two
specified tasks. At smaller order twelve, the preceding one-case synthesis
is unavailable:
no source theorem handles almost every parameter. We instead prove an
explicit eventual threshold by an elementary coefficient estimate, then
certify the entire remaining finite set over one finite field. The proof
uses only six-square retained matrices, even when the original matrix has
more than eleven thousand rows.

\begin{theorem}[Complete smaller-order-twelve reconstruction]
\label{t12:complete}
For every integer $b\ge7$, the matrix $M_{6,b}$ defined in
\eqref{t10:M} is nonsingular. Consequently, for every even $q>12$,
the first $q$ centered moments of $A_{12}(n)A_q(n)$ have formal
coordinate rank $q$, and every later centered moment is a finite
rational linear combination of them.
\end{theorem}

The proof has two exact components:
\[
 \det M_{6,b}\ne0\quad(b\ge11520),\qquad
 \det M_{6,b}\ne0\quad(7\le b\le11519).
\]
The first component has an ordinary analytic proof below. The second is
a finite arithmetic certificate with $11513$ cases, reproduced in full
by two independently implemented verifiers. No sampling of the parameter
range is used.

\subsection{A six-square polynomial and its leading determinant}

In \eqref{t10:Newton}--\eqref{t10:A}, set $a=6$ and replace
$2B_{2b-2j}$ by independent indeterminates $U_j$, $1\le j\le11$.
Write the resulting six-square matrix as $\mathcal A(b;U)$. Set
\begin{align*}
 n_i&=b-6+i,&R_\ell&=2b-10+2\ell,\\
 J(b)&=\prod_{i=0}^5(2n_i+1)(2n_i+2),&
 E(b)&=\prod_{i,\ell=0}^5
       \prod_{r=1}^{\max(0,2i-4\ell-1)}(R_\ell+r),\\
 D(b)&=12!J(b)E(b),&
 \mathscr P(b;U)&=\det(D(b)\mathcal A(b;U)).
\end{align*}
The polynomial $E$ has $35$ linear factors. The source factorial
cancellation proves
\[
 D\mathcal A\in\mathbb Z[b,U_1,\ldots,U_{11}]^{6\times6},
 \qquad\deg_b(D\mathcal A_{i\ell})\le67,
 \qquad\deg_b\mathscr P\le402.
\]
Indeed $N_{i\ell}-R_\ell=2i-4\ell-1$, so $E$ clears exactly the
possible denominator products in $R_\ell!/N_{i\ell}!$;
$J$ clears the row factors and $12!$ clears the fixed factorials.
All these factors are positive when $b\ge12$.

Let $\mathscr P_{[6]}$ denote the part homogeneous of degree six in
the $U$ variables. The inherited highest-degree determinant identity,
specialized to six, gives at $U_j=2B_{2b-2j}$
\begin{equation}\label{t12:leading}
 \left|\mathscr P_{[6]}\right|
 =\frac{\mathcal G(b)(4F_0)^6}{(2\pi)^{12(b-6)}}
   \det\bigl(\zeta(s_0+2i+2j)\bigr)_{0\le i,j<6},
 \qquad s_0=2b-22,\quad F_0=s_0!,
\end{equation}
where
\begin{align*}
 \mathcal G(b)={}&\kappa J(b)^5E(b)^6
 \prod_{i=0}^5(4n_i+2)^{\underline{2i}}
 \prod_{\ell=0}^5\prod_{r=1}^{12+2\ell}(2b-22+r),\\
 \kappa={}&\frac{6!(12!)^6}{\prod_{\ell=0}^5(2\ell+2)!}
 =3590542128270039916216320000000.
\end{align*}
The degree of $\mathcal G$ is $402$, and every linear factor displayed
in it is at least $b$ when $b\ge24$. Hence
\begin{equation}\label{t12:Gbound}
 \mathcal G(b)\ge\kappa b^{402}\qquad(b\ge24).
\end{equation}

Here the leading determinant is a structural identity, not an estimate
of an expanded polynomial. Its short derivation uses
$N_{i\ell}-(2k+1)=2b-2(6-i+\ell+k)$ in \eqref{t10:A} to cancel
factorials. The resulting linear-in-$U$ part factors through the
matrix $(V_{6-i+\ell})$, where
$V_j=U_jF_0/(2b-2j)!$. The remaining row transformation is triangular,
with diagonal
$d_{i1}=-(i+1)(4n_i+2)^{\underline{2i}}$.
Taking determinants, reversing rows, and applying Euler's formula
$2B_{2m}=(-1)^{m-1}4(2m)!\zeta(2m)/(2\pi)^{2m}$ gives
\eqref{t12:leading} with the displayed $\mathcal G$. This retains the
attribution of the general highest-degree factorization in the source
\emph{Quartic Spectral Reconstruction}.

The zeta determinant in \eqref{t12:leading} is a positive moment
Gram determinant: use the points $r^{-2}$ with weights $r^{-s_0}$,
$r\ge1$. Cauchy--Binet and the first six distinct support points give
\begin{equation}\label{t12:zeta-lower}
 \begin{aligned}
 \det\bigl(\zeta(s_0+2i+2j)\bigr)_{0\le i,j<6}
 &\ge V\,720^{-s_0},\\
 V&=\prod_{1\le r<s\le6}(r^{-2}-s^{-2})^2\\
 &=\frac{14235529}{3743906242624487424000000000000}>0.
 \end{aligned}
\end{equation}
In particular, the highest-degree component cannot vanish.

\subsection{A simple coefficient estimate requiring no expanded determinant}

For an integer polynomial let $\|\cdot\|_1$ be the sum of absolute
values of its coefficients. It is submultiplicative. We shall use
the deliberately generous bound
\begin{equation}\label{t12:norm}
 \max_{0\le i,\ell<6}\|D\mathcal A_{i\ell}\|_1<2^{512}.
\end{equation}
The following elementary factor count proves it.
Every factor in $J$ has norm less than $16$, and every factor in $E$
has norm less than $16$. Thus
\[
 \|J\|_1<2^{48},\qquad\|E\|_1<2^{140},\qquad12!<2^{29}.
\]
After removing the two factors for one fixed row, the remaining
$J$ product has norm less than $2^{40}$. After cancellation in
$ER_\ell!/N_{i\ell}!$, at most $21$ additional linear factors are
present, each of norm less than $32$. Its norm is therefore less
than $2^{245}$. Also
\[
 \|N_{i\ell}^{\underline h}\|_1<2^{55}
 \quad(1\le h\le11),\qquad
 \|d_{ih}\|_1<2^{65}.
\]
For the last assertion, apply the Newton formula to
\[
 W_i(j)=(2b-10+2i-j)
             \prod_{r=0}^{2i-1}(4b-22+4i-j-r),
 \qquad0\le j\le h\le2i+1\le11.
\]
Its first factor has norm below $16$ and each remaining factor
has norm below $32$, so $\|W_i(j)\|_1<2^{54}$. The sum of the
binomial coefficients in the Newton formula is $2^h\le2^{11}$;
division by $h!$ only decreases the bound.

The numerator defining $F_{i\ell}$ has at most $20$ linear factors,
each of norm below $32$. Thus the constant term $2DF_{i\ell}$ has
norm below $2^{318}$. Each coefficient of an indeterminate $U_j$
in $D\mathcal A_{i\ell}$ has norm below
\[
 2^{29+40+245+65+55}=2^{434}.
\]
There are at most six such coefficients. Consequently the entry norm
is below $2^{318}+6\cdot2^{434}<2^{437}<2^{512}$, proving
\eqref{t12:norm}. This argument also proves integrality by the same
explicit cancellations; no fitted coefficient bound is involved.
A supplementary exact script also computes a sharper factor-product
bound; the proof uses only the elementary bound just established.

The determinant coefficient norm is at most
\[
 K=6!(2^{512})^6=720\,2^{3072}.
\]
For $b\ge11520$ put
\[
 T=\frac{8(2b)^{20}F_0}{(2\pi)^{s_0}}>1.
\]
The inequality $T>1$ follows from $s_0!/(2\pi)^{s_0}>(s_0/24)^{s_0}>1$,
using $s_0!>(s_0/3)^{s_0}$ and $2\pi<8$.
Euler's Bernoulli formula and $\zeta(2m)<2$ show
$|2B_{2b-2j}|\le T$ for $1\le j\le11$.
Since all remaining terms of $\mathscr P$ have total $U$ degree
at most five,
\[
 |\mathscr P_{[\le5]}|\le Kb^{402}T^5.
\]
Combining this estimate with \eqref{t12:leading}--\eqref{t12:zeta-lower},
then using $2\pi<8$, gives the explicit domination ratio
\begin{equation}\label{t12:ratio}
 \frac{|\mathscr P_{[\le5]}|}{|\mathscr P_{[6]}|}
 \le C\,\frac{b^{100}5760^{2b}}{(2b-22)!},\qquad
 C=\frac{2^{3289}}{\kappa V\,720^{21}}<2^{3289}.
\end{equation}
The final inequality follows, for example, from the exact rational value
\[
 \kappa V=\frac{786377640075797}{57600000}>1.
\]

\subsection{A threshold of 11520 proved by elementary inequalities}

Define
\[
 R(b)=2^{3289}\frac{b^{100}5760^{2b}}{(2b-22)!}.
\]
We prove $R(b)<1$ for every integer $b\ge11520$. At $b=11520$,
write $n=23018$. The elementary factorial estimate
$n!\ge(n/e)^n>(n/3)^n$, obtained by integrating $\log x$, gives
\[
 n!>7200^n=5760^n(5/4)^n,
\]
since $n>3\cdot7200$. Moreover $11520<2^{14}$,
$5760<2^{13}$, and $(5/4)^4=625/256>2$. Therefore
\begin{equation}\label{t12:initial-bound}
 R(11520)
 <2^{3289+1400+286}(4/5)^{23018}
 <2^{4975-5754}=2^{-779}<1.
\end{equation}
For every subsequent integer $b$,
\[
 \frac{R(b+1)}{R(b)}
 =\left(1+\frac1b\right)^{100}
       \frac{5760^2}{(2b-20)(2b-21)}<\frac34,
\]
because $b\ge100$ gives the first factor less than $3$,
both denominator factors exceed $b$, and $b\ge11520=2\cdot5760$.
Thus the ratio in \eqref{t12:ratio} is strictly below one for all
$b\ge11520$. Its highest-degree component cannot be cancelled.
As $D(b)>0$ and the source multiplier \eqref{t10:multiplier} is
nonzero, this proves $\det M_{6,b}\ne0$ throughout that infinite
range. The explicit threshold follows from a mathematical argument,
not from the observed behavior of a finite list of determinants.

\subsection{The complete finite-field certificate below the threshold}

For the remaining $11513$ integers $7\le b\le11519$, use the
single prime $P=65537$. The original source matrix is used for
$7\le b\le11$, where the contraction is not asserted. For
$12\le b\le11519$, use $\mathcal A$ from \eqref{t10:A} with $a=6$.
The source relation \eqref{t10:multiplier} implies that nonvanishing
of this six-square determinant suffices.

Every division is defined modulo $P$. The largest Bernoulli index
needed for the retained matrices is $2\cdot11519-2=23036<P-1$;
\eqref{t10:Brec} therefore puts all these Bernoulli numbers in
$\mathbb Z_{(P)}$. Every factorial argument and every positive
linear denominator appearing in the matrices and the exact source
multipliers is less than $4\cdot11520+5=46085<P$. In particular,
the entire source multiplier is a unit modulo $P$. The computation
never replaces a singular reduction modulo $P$ by an unsupported
claim of rational nonvanishing: every recorded residue is nonzero.

The file \texttt{complete\_p12.json} records the full ordered residue
list, with its first entry corresponding to $b=7$. It has $11513$
entries, and every entry belongs to $\{1,\ldots,65536\}$. For orientation,
a few entries are
\begin{equation}\label{t12:sample}
\begin{array}{c|rrrrrrrr}
 b&7&8&11&12&100&962&4096&11519\\\hline
\text{residue}&39450&1561&25713&60062&43034&42605&25983&23212.
\end{array}
\end{equation}
For $b\le11$ the entry is $\det M_{6,b}\bmod P$;
for $b\ge12$ it is $\det\mathcal A\bmod P$.
The displayed sample is explanatory; the proof checks the complete
ordered list.

Two full implementations replay the certificate by different exact
algorithms. The C++ verifier uses the triangular Bernoulli recurrence,
finite differences for the Newton coefficients at each integer $b$,
and modular Gaussian elimination. The Python verifier instead
inverts $(e^t-1)/t$ as a truncated power series, using exact integer
Kronecker multiplication. This produces $B_n/n!$ directly. It then
uses the factorial-cancelled entry formula
\begin{equation}\label{t12:cancelled-entry}
 \mathcal A_{i\ell}=2F_{i\ell}
 +\frac{R_\ell!}{(2n_i+1)(2n_i+2)(2\ell+2)!}
 \sum_{k=0}^i d_{i,2k+1}
 \frac{2B_{2b-2(6-i+\ell+k)}}{(2b-2(6-i+\ell+k))!},
\end{equation}
with Newton coefficients computed as global integer polynomials
in $b$, and evaluates determinants by a division-free subset
Laplace expansion. It verifies every one of the $11513$ residues,
checks that they are all nonzero, and compares the complete result
with the independently produced C++ list. The C++ calculation also
compares the original and contracted determinant formulas throughout
$12\le b\le32$, including the full exact multiplier.

The exact polynomial multiplication in the Python verifier represents
coefficients as successive base-$2^{64}$ digits of an integer. At the
largest product length here, every unreduced convolution coefficient
is smaller than $2^{64}$; hence there is no carry between coefficient
blocks. Integer multiplication is therefore precisely polynomial
convolution, with no floating-point transform or rounding. Newton
inversion is checked by multiplying the final series by
$(e^t-1)/t$ through degree $23039$ and obtaining the constant polynomial
one. Both complete replays use finite integer arithmetic. Together
with \eqref{t12:initial-bound}, they prove
Theorem~\ref{t12:complete}.

\subsection{Identity consequences and the remaining problem}

The new rank theorem yields the ordinary convergent identity family
\eqref{t10:ordinary} with the following replacements:
$M_{5,b}$ by $M_{6,b}$, $10$ by $12$, $9!$ by $11!$, and
$\psi^{(9)}$ by $\psi^{(11)}$. In the subtraction polynomial,
replace the lower outer index $b+4$ by $b+5$ and each occurrence of
$r-b-4$ by $r-b-5$. The allowed range is $b\ge7$. Thus products
of the eleventh polygamma with every higher odd-order polygamma
now reconstruct all $b$ even-zeta coordinates by explicit rational
polynomial weights and finite midpoint subtractions.

Theorems~\ref{t10:complete} and \ref{t12:complete} extend the proved
smaller orders from $2,4,6,8$ to $10,12$. The full fixed-gap-six
family and the unrestricted unequal-order conjecture remain open.
The eventual theorem reduces each further fixed smaller order to a
finite certificate, but no assertion about uncomputed orders is made
here. The leading Hankel factorization and the source contraction
are inherited; the completed order-twelve family, its simple explicit
threshold proof, and its full pair of arithmetic certificates are the
new contributions of this section.
